\subsection*{把发送放在哪里}

承接上一章 P1/C1/T1 的恢复链，增加一项通知履约系统的责任。为比较发送次序，以下每条历史都从购买库、平台收款和队列为空开始：平台账户初始 10000 分，Q7 对应 CNY 3000 分。A 保存原购买及固定平台身份，E 建立 X1、账户变为 7000 分，R 建立 O1、保存原结果，令 P1 为 \code{paid}、T1 为 \code{done}。完整购买身份、商品和输入解释沿用第七章。

现在增加一条消息，告诉履约系统订单 O1 已付款。消息系统中的代理负责接收发布，并向订阅者分发消息。购买库的 R 与代理的接收分别有自己的提交决定。购买服务要修改数据库，又要把通知写进消息系统，这种跨两个保存位置的写入通常称为双写。这里的两次写入有共同的业务目的，却各自决定是否保存成功。数据库的提交确认对应订单与结果，代理的受理确认对应消息。恢复者分别核验这两份证据，再决定哪一段工作要继续~\cite{outbox}。

先看提交 R 以后再发。R 的提交确认已经到达，进程尚未发送就退出。数据库恢复后保留 O1 和原结果，代理中没有消息，内存中准备发送的变量也已经消失。重复查询 P1 可以恢复购买结果，但查询路径只返回结果，发布所需的工作已经随进程消失。恢复者要继续通知，数据库就得另外留下可查的发布依据。

再把发送放到 R 提交以前。代理先接收并交付消息，履约系统开始处理 O1。随后 R 因约束异常回滚，数据库里没有已提交的 O1。代理独立保留消息，消费者也可能已经保存自己的处理结果。让消费者等待一会儿再查询，只是改变间隔，R 可以在等待以后继续阻塞或回滚。

在 R 的事务内部调用代理，也会出现这两种先后次序。SQL 事务的开启与结束限制数据库写入，网络调用一旦交给另一个系统，对方即可独立提交。代理先接收、数据库后回滚，会留下没有对应订单的通知。数据库先提交、代理后接收，进程则可能在两者之间退出~\cite{outbox}。

\CNineFigure{ch09-original01}{两种直接双写次序分别留下漏通知与先通知的历史}{fig:ch09-dual-write}

这个问题与第七章的外部付款相似，不过这里可以改变通知的表达。付款效果必须由平台执行，待发送消息的全部语义却可以先保存在购买数据库中。第一次事务的目标由此变为：业务成功时，同时留下一项以后必须继续发布的工作。

\subsection*{先保存待发送的事实}

在 R 内插入一行事件记录，写明 O1 已付款以及发布所需的内容，再与 O1、原结果和 T1 完成共同提交。这张表称为 Outbox，可以理解为与业务数据放在一起的待发事件表。随后由独立的发布器扫描这些记录，向代理发送~\cite{outbox}。

R 在提交以前退出，订单和待发事件一起回滚。R 已经提交，进程即使立即退出，后来发布器也能查到事件。数据库的原子性连接业务事实与发布依据，代理则在稍后的独立提交域中受理消息。

事件与消息副本由此分开。一项事件表达一次已保存的业务事实，发布器可以为它发出多次网络尝试，代理也可能因此保留多个消息副本。R 创建事件的次数、发布函数被调用的次数，以及代理接收副本的次数，分别描述不同的动作。

Transactional Outbox 原模式已经指出，转发进程在代理接收以后、本地记录以前退出，会造成重复发布。\footnote{原模式对转发进程写道“might publish a message more than once”。这里的重复指同一业务事件可能再次发布，事件唯一与代理副本唯一分别由各自边界解释~\cite{outbox}。}接下来的设计要保住这项事件，并让重发的每个副本携带同一个事件身份。下游怎样把重复副本归并为一次业务效果，将在第十章推导。

下面先从共同空初态推导新 R：事件保存什么，怎样与成功共同提交，以及发布者如何继续工作。第八章已经成功的旧 O1 没有这项事件，升级时要补建其通知。第~\ref{section:10-3}~节在消费者机制建立以后，完整处理源端补建与旧履约迁移的两道屏障。
